\chapter[Sermon 6. How John Was Affected Towards the Vision to Him Exhibited, the Comfort of John, and the Exposition of the Vision, Applied unto Consolation.]{How John Was Affected Towards the Vision to Him Exhibited, the Comfort of John, and the Exposition of the Vision, Applied unto Consolation.}
\label{sermon:006}
\markright{Sermon 6. How John Was Affected Towards the Vision to Him ...}

And when I saw him, I fell at his feet as if dead. And he laid his right hand upon me, saying to me, “Fear not. I am the first and the last, and I am alive, and was dead. And behold, I am alive forevermore, and have the keys of hell and of death. Write therefore the things which you have seen, and the things which are, and the things which shall be fulfilled hereafter. And the mystery of the seven stars which you saw in my right hand, and the seven golden candlesticks. The seven stars are the messengers of the seven congregations. And the seven candlesticks which you saw are the seven congregations.

It follows how blessed John was moved by that heavenly and wondrous vision, and how he received comfort, along with an explanation of the vision applied to his comfort, with a command to write down all these things diligently.

\index{Daniel (Book of)}What time he had fully seen this divine and heavenly vision of Christ our Lord, sitting on the right hand of God in glory, his strength failing him, he falls down on the earth: and lying at the feet of the Lord, is like a dead body. We read that the same happened to Daniel in the tenth chapter. And other men of God also have been frightened by the visions of Angels. The women also in the New Testament trembled at the sepulchre, seeing Angels. Peter was amazed at the greatness of the miracle. Luke 5. And falling at the knees of the Lord, cries out, “Go from me, Lord, for I am a sinful man.” For godly visions betray our infirmity. Neither are we prepared or sufficiently purged to behold those celestial things with eyes and minds sick and not yet well purified. Therefore, the elect must be glorified in another life, that they may be made partakers of celestial glory. In the mean season here, all the godly are humbled and abased by high visions and revelations. For they do not advance themselves proudly through the glory of revelation, but perceiving their natural corruption, they crave pardon and the augmentation of the celestial grace and light. For unless we be illumined with the spirit of God, we shall lie like dead people, however we receive with our corporal senses the visions of heaven.

But those who humble themselves before the Lord find a most present consolation at the Lord’s hand. Therefore, there came to John immediately both in word and in a full consolation. For the angel representing the person of Christ lays his right hand upon John, which is a token of friendship, protection, and of present help. For when pressing this manner of speaking in Dutch, we say that laying on the hand signifies that Christ is good to John, ready to help him. Which he immediately makes plain by the addition of words, saying, “Fear not.” This saying is common everywhere in the story of the Gospel, and therefore is most gospel-like, that is to say, most fortunate. For God commands the humbled to be of good hope and to live assured under the protection of the highest. Which truly we understand to be spoken not to John alone but to all of us also, that we in like manner, albeit that we feel the infirmity of our flesh, should yet hope well of the goodness and mercy of God. Here follows that cause more fully declared, why John should not be afraid. For the vision showed was not exhibited for the terror of him, but that John might perceive how great and mighty he is who is prepared for the defense of him and all the faithful. As though he should say, “Where you see how great he is, who has taken upon himself to defend you, who finally protects and governs the whole Church, there is no cause why you should be afraid. But rather execute boldly what he commands you. Write what he commands to be written. Be not afraid of men, fear God rather. For if good men are so sore afraid at the sight of him, where shall the enemies and contemners of God appear?”

Therefore, consequently, he explains the vision, teaching who he is, who was seen like the Son of Man walking among the golden candlesticks. And he applies this explanation to comfort, so that both John and every faithful person may perceive how mighty Christ is, and what the faithful have obtained through him. For the angel shapes his speech so that we may seem to hear all things spoken not by the mouth of the angel, but by Christ himself. And this explanation has its parts. First, he declares (as I said just now) whose image it was that was shown. Then, a commandment to write this book is added. After that, the mystery of the stars is opened. Finally, the secret of the candlesticks is revealed: and all these things, right plainly and briefly.

First you have seen, says the Lord, a vision whereat you were dismayed: but fear not. For you have not seen any evil or fearful spirit, foretelling any misfortune: but my form, which am your redeemer and Lord. I am first and last. And this manner of speaking he took out of the Prophecies of Isaiah, as it is to be seen in the 41st, 43rd, 45th, and 48th chapters. And he signifies himself to be coequal and of the same substance with the Father in all things, very God, eternal, and incomprehensible. For just as the Father attributes things to himself, the Son also claims them. But there is no order or certain time to be understood in “first and last,” but plainly everlastingness. Therefore Christ here signifies that he is very God, equal and of the same essence with the Father from all eternity, as is also much confirmed in John 1:5, 10, 14, and 17. This contends against the heretics, who at that time also, as today the Servetarians, deny the eternal deity of Christ the Lord. And thus, when the true God is acknowledged and believed by us, he may be for our salvation. If Christ is not very God, he is not our salvation. For I am God, says the truth: And besides me there is no God, no salvation.

Secondly, he says, “I am alive, and was dead,” by which he signifies that he took on the true human nature. Which many also at the same time denied, just as there are some at this day who plainly diminish the humanity of Christ. Against all such heresies, the Lord himself confesses that he was dead. Whereby it is now manifest that he is very man, as he is also very God, of the same essence with his Father in deity, as he is also of the same substance with us in humanity, like unto us in all things, sin excepted. For he took not the nature of angels, but the seed of Abraham. And it behooved indeed that the Son of Man should be incarnate, that both he might suffer and shed blood. For the testament in the dead is finally ratified; neither is there any remission made without shedding of blood. The Lord therefore dies and sheds blood, to the intent he might give full remission of sins and confirm the new testament. Yet even he that was thought to be dead, now lives, and is that same living one, who having vanquished death, on the third day rose again from the dead, and restores life to all believers, and inspires into them his own very life.

And therefore adds immediately: behold I am living world without end. For now Christ dies no more, death shall not rule over him. But rather he is the life of all his faithful, who in rising again, brought again life: and that life everlasting, enduring, I say, world without end. As he himself declares more at large in John, chapter 5, verses 6 and 10. And the Apostle to the Romans, 4:1; Corinthians, 15; and 2 Timothy, 1.

Furthermore, where many were inclined to doubt this life obtained and restored by Christ, the Lord himself confirms what he said through another, and says, “Amen.” As though he should say, “This is altogether true and certain that I say.”

\index{Antichrist}\index{Papacy}\index{Resurrection}\index{Rome}Finally, he adds, “And I have the keys of Hell and of death.” By these words, he greatly comforts and expresses his power, declaring how great he is and what we have from him. Here, we must speak by way of analogy. The common explanation says well: he who has the keys to any house lets in whom he will and keeps back whom he will from entering. Therefore, Christ possesses the keys of death and Hell, for he delivers whom he will from perpetual condemnation of death, and whom he will, he allows to remain justly in the same danger of damnation. Indeed, Isaiah in chapter 22, speaking of Eliakim, whom he says should be made steward in the court of Hezekiah, says, “I will lay the key of the house of David upon his shoulder, which shall open, and no man shall shut; shall shut, and no man shall open.” Therefore, keys in Scripture represent the charge and governance of a house. Eliakim would govern all things in the court of Hezekiah uprightly. Whatever he determined, no one would infringe; what he abrogated, no one would restore. Christ, of whom Eliakim was a figure, himself also has the chief governance in the house or kingdom of God, so that he may quicken whom he wills and pluck back from Hell and damnation. Again, he may condemn whom he wills through his just judgment. For he has full power over death and Hell. For he has overcome and weakened both. These things greatly comfort the faithful and retain them in all godly duties. It is chiefly to be observed that he does not say he \textit{had} the keys or \textit{will have} them, but “I have,” he says, “I have.” He did not give his power to the Bishop of Rome, but retains it himself and will keep it forever. Nor did he give to the Apostles full power of life and death, of salvation and damnation, thus divesting himself; but he gave the keys of opening and shutting heaven, as it were, to his ministers and servants, through the preaching of the Gospel. By this, he promised life to all who believed, and Christ himself should give that life for the truth of the promise. To whomever they should threaten damnation, Christ himself should condemn for the truth of his word. Therefore, we see that the Lord still keeps and exercises the power, and his ministers the ministry (by preaching, not by absolute power). Therefore, the Pope is Antichrist, who usurps and takes upon himself this full power and authority in Heaven and on Earth, and in the midst of the earth, or beyond all the earth, namely those unfortunate islands. I mean purgatory. By this crafty device, he has subtly emptied the purses, coffers, garners, and wine cellars, of foolish people who stray from the articles of their belief, namely, that I believe in the forgiveness of sins, the resurrection of the flesh, and everlasting life. The beast dares usurp the two horns of the Lamb, namely, the authority of King and Bishop, and therefore hangs two keys[illegible] under his triple crown, so that even by these emblems all the world may perceive that this is he who, having subdued three kings or horns, has arisen and claims all power in heaven and on earth, signified by the two keys. And surely the blindness of our time is wonderful and to be lamented, that having eyes, it sees nothing. Let those who are wise remember that Christ still has the keys of death[illegible] and Hell, his ministers the denouncing of life and death.

And now when he had declared these great and wholesome matters, and had comforted the mind of John, he adds the commandment: write the vision exhibited, finally write those things also which must be done shortly after this. He places in the midst those that be—that is, which are in deed, and true, and be not false. And these things are to give authority to this book, finally to the whole scripture, which is revealed with like truth of the same Author. And as John is commanded to write without fear, so are we commanded to preach and publish the same boldly, though the world be never so mad thereat.

He adds further, the exposition that remains, and says: “The mystery of the seven stars, \&c.” The reason seems almost unpersisted [?]. Therefore, we must understand this to be the mystery or sacrament of the seven stars and candlesticks, so that it may be as it were a proposition, and the exposition should follow immediately. The seven stars are seven messengers, \&c. And by sacrament understand a secret mystery, and the very exposition of the mystery. As if to note here the goodness of Almighty God, which declares to us even the hardest places of Scripture. Where, then, are those who accuse the Scripture of obscurity, and contend that it cannot be understood? Let us also mark here the common manner of speaking of the whole scripture: seven stars are seven messengers. The seven lights are seven churches. For signs receive the names of the things they are, although they remain in their own substance, and are not changed into another. This is also granted by those contentious persons who, in the words of the supper, “This is my body,” will acknowledge no figurative speech at all.

Stars are called angels. Angels are God’s messengers, pastors of churches, so called in the second and third chapters of Malachi. For God sends preachers as ambassadors to the people, and wills them to be heard in the same way as himself. Luke 10; John 13. Let no man therefore wait until the Lord himself comes down from heaven again and preach to us. Even now he preaches to us by his messengers, who preach his word, that is to say, the word of Christ. If you despise them, you despise Christ. Preachers are called stars because of their bright and heavenly doctrine, and for their purity of life. Beware therefore you preachers, that you be not wandering plants, lest you have no light at all, neither in doctrine nor in conversation of life. For then you shall be likened to stars that fall out of heaven, as will happen hereafter in this book to false teachers.

But those stars are not in the head, or in the feast, or on the back or sides: but in the right hand of Christ. This thing has indeed a great consolation for the pastors, being in God’s right hand, in God’s protection, so that no man shall take them out of his hand. God himself also upholds pastors, and furnishes them with necessary goods of the Church. Therefore is the whole government and glory his. Wherefore the Apostle says also: he that waters and plants is nothing, but God that gives increase.

\index{Arethas of Caesarea}Now as concerning the candlesticks, there was certainly in the Tabernacle of Moses with seven sockets, set in seven candles. In Solomon’s temple there were ten candlesticks. The one represented a figure of Christ: and even upon it, and the ten, signified the universality of churches, which are lit by the only light Christ, and have of this one, whatever light they have. And the candlesticks are of gold. The mystery whereof Arethas expounds: They are all gold, says he, for the purity and preciousness of faith lying hid in them. Indeed, the candlesticks of themselves give no light, but are receptacles of light. So from us arises no light, but darkness. But in that everlasting light, set a light in the candlestick, the light shines: if Christ illumine the church with faith and understanding, then faith shows itself in open confession and the purity of life in conversation. And this the Lord requires of his church in the fifth chapter of Matthew. So let your light shine, and so forth. And the apostle in the second letter to the Philippians: Shine like lights in the world amidst a crooked and perverse generation.

And hitherto we have handled the consolation of Christ and the exposition of that great and celestial vision, whereby we have learned the mysteries of the faith of Christ and of his Church: to the end we should know that Christ is the Lord reigning in his Church, and applying all things to the salvation of his faithful. That he sends preachers, teaches by them, and keeps and defends them. To him be glory, and so forth.
